\chapter{Introdução}\label{cap:introducao}

Os \glspl{LLM} ampliaram as possibilidades de automação de atividades que envolvem compreensão e geração de linguagem natural. Quando integrados a bases de conhecimento, ferramentas e mecanismos de decisão, esses modelos podem formar agentes capazes de interpretar solicitações, manter o contexto de uma conversa e executar tarefas com algum grau de autonomia. No atendimento ao cliente, essa abordagem permite desenvolver soluções mais flexíveis do que os chatbots baseados apenas em regras e fluxos predefinidos.

Trabalhos recentes apresentam aplicações de LLMs na classificação de chamados e na geração de respostas de suporte~\cite{moncada2025}, bem como arquiteturas formadas por agentes coordenadores e especializados para atuar em diferentes componentes de sistemas de telecomunicações~\cite{lee2025}. Essas soluções podem ser organizadas em arquiteturas de agente único ou multiagente.

Na arquitetura de agente único, um mesmo componente interpreta as mensagens, consulta as informações disponíveis, decide o procedimento e produz as respostas. Essa centralização reduz a necessidade de coordenação entre componentes. Na arquitetura multiagente, essas responsabilidades são distribuídas entre agentes com funções específicas, como triagem, primeiro atendimento, comercial, ouvidoria e suporte técnico. A divisão pode facilitar a especialização e reduzir o contexto recebido por cada agente, mas também introduz etapas adicionais de comunicação e novos pontos de falha~\cite{yan2026}.

A comunicação entre os agentes influencia a eficiência, a escalabilidade e a robustez do sistema. Uma decisão incorreta na triagem pode encaminhar a conversa para um agente que não possui as informações necessárias. Também podem ocorrer perdas de contexto, encaminhamentos desnecessários ou propagação de erros entre os componentes. Falhas de agentes e interrupções na comunicação também são apontadas como riscos desses sistemas~\cite{owotogbe2025}.

Estudos recentes investigam o encaminhamento de conversas para atendimento humano e o reconhecimento de intenções em ambientes multiagentes~\cite{amri2026,yang2025}. Os resultados apontam benefícios da especialização, mas também mostram que a vantagem da arquitetura multiagente não ocorre em todas as métricas. Além disso, os trabalhos adotam diferentes modelos, cenários e critérios de avaliação, o que dificulta uma comparação direta entre arquiteturas de agente único e multiagente. Na comparação realizada por Amri et al.~\cite{amri2026}, também foi observada uma tendência de o sistema multiagente realizar transferências para atendimento humano mais cedo.

Esses resultados indicam que a adoção de uma arquitetura multiagente não representa, necessariamente, uma melhoria em todas as situações. Ainda são necessárias comparações nas quais as arquiteturas utilizem o mesmo modelo, a mesma base de conhecimento, parâmetros equivalentes e cenários de atendimento comparáveis. Também é necessário observar o desenvolvimento da conversa, as decisões de encaminhamento e os recursos utilizados até a conclusão da tarefa.

\section{Problema de pesquisa}\label{sec:problema}

A escolha entre uma arquitetura de agente único e uma arquitetura multiagente envolve um compromisso entre qualidade, complexidade e custo operacional. A especialização pode melhorar o tratamento de determinadas solicitações, mas exige coordenação entre agentes e pode aumentar o número de chamadas ao modelo e o consumo de tokens.

Uma avaliação baseada somente em perguntas isoladas não representa todo o funcionamento de um chatbot de atendimento. Em uma conversa real, as informações são fornecidas gradualmente, e uma resposta pode depender das mensagens anteriores. O sistema também pode precisar rever o encaminhamento quando o usuário apresenta uma nova informação ou muda de assunto durante o atendimento.

Diante desse contexto, este trabalho busca responder à seguinte pergunta:

\begin{quote}
Quais diferenças de desempenho são observadas entre uma arquitetura de agente único e uma arquitetura multiagente, ambas baseadas em modelos de linguagem em grande escala e aplicadas a cenários conversacionais de atendimento, quanto à correção e à relevância das respostas, à conclusão das tarefas, à precisão dos encaminhamentos, ao número de chamadas ao modelo e ao consumo de tokens?
\end{quote}

\section{Objetivos}\label{sec:objetivos}

\subsection{Objetivo geral}\label{subsec:objetivo-geral}

Comparar experimentalmente o desempenho de uma arquitetura de agente único e de uma arquitetura multiagente, ambas baseadas em modelos de linguagem em grande escala e aplicadas a cenários conversacionais de atendimento textual de um provedor de Internet.

\subsection{Objetivos específicos}\label{subsec:objetivos-especificos}

Para alcançar o objetivo geral, foram definidos os seguintes objetivos específicos:

\begin{itemize}
    \item definir uma base de conhecimento e regras de atendimento comuns às duas arquiteturas;

    \item especificar e implementar um protótipo de agente único e um protótipo multiagente composto por um agente de triagem e agentes especializados;

    \item elaborar e revisar um conjunto de cenários conversacionais de teste para primeiro atendimento, comercial, ouvidoria e suporte técnico;

    \item avaliar o desempenho das duas arquiteturas quanto à correção e à relevância das respostas, à conclusão das tarefas, à precisão dos encaminhamentos, ao consumo de tokens, ao número de chamadas ao modelo e ao custo estimado de execução;

    \item comparar os resultados para identificar as vantagens, limitações e custos de cada arquitetura.
\end{itemize}

\section{Delimitação do trabalho}\label{sec:delimitacao}

O estudo será restrito ao atendimento textual e utilizará
cenários sintéticos baseados em situações típicas de atendimento
de um provedor de Internet. Todos os dados empregados serão
fictícios, sem utilização de conversas ou cadastros de clientes
reais e sem interação com clientes durante o experimento.

As duas arquiteturas utilizarão o mesmo modelo de linguagem, os mesmos parâmetros de geração e a mesma base de conhecimento. No agente único, as regras e informações serão reunidas em um único contexto. No sistema multiagente, o conteúdo será distribuído conforme as responsabilidades da triagem e dos agentes especializados.

Serão consideradas quatro áreas de atendimento: primeiro atendimento, comercial, ouvidoria e suporte técnico. Os cenários poderão conter mudanças de assunto durante a conversa, permitindo verificar o comportamento das arquiteturas diante da necessidade de um novo encaminhamento.

Não fazem parte do escopo o atendimento por voz, o processamento de conteúdo multimodal, o treinamento ou ajuste fino de modelos, a comparação entre diferentes LLMs e a implantação do sistema com clientes reais. Também não serão realizados testes de robustez com injeção deliberada de falhas nos agentes ou na comunicação.

A latência das respostas não será avaliada. Como o acesso ao modelo será realizado por uma API externa, o tempo observado poderá variar conforme as condições da rede e do serviço, sobre as quais o experimento não terá controle. A comparação dos recursos utilizados ficará restrita ao número de chamadas ao modelo, ao consumo de tokens e ao custo estimado de execução.

\section{Organização do texto}\label{sec:organizacao}

Além deste capítulo introdutório, a proposta está organizada da seguinte forma: o \autoref{cap:revisao} apresenta a fundamentação teórica, os trabalhos relacionados e a lacuna de pesquisa;
o \autoref{cap:metodologia} descreve as arquiteturas propostas, a base de conhecimento, os cenários conversacionais, as ferramentas, o procedimento experimental e os critérios de avaliação;
e o \autoref{cap:cronograma} apresenta o cronograma das atividades previstas para o desenvolvimento do trabalho no TCC2.